%----------------------------------------------------------------------------------------
%	IMPERIAL COLOUR PALETTE
%	xcolor is already loaded by MastersDoctoralThesis.cls (with dvipsnames).
%	Names below match Imperial's supplied swatches; reuse them in tables, tikz,
%	and (mirrored in Python) for figures.
%----------------------------------------------------------------------------------------

% ---- Primary ----
\definecolor{ImperialBlue}{HTML}{0000CD}   % headline colour
\definecolor{ImperialNavy}{HTML}{000080}
\definecolor{ImperialDark}{HTML}{232333}

% ---- Extended palette ----
\definecolor{ImpSaddleBrown}{HTML}{8B4513}
\definecolor{ImpTeal}{HTML}{008080}
\definecolor{ImpVioletRed}{HTML}{C71585}
\definecolor{ImpIndigo}{HTML}{4B0082}
\definecolor{ImpCrimson}{HTML}{DC143C}
\definecolor{ImpOrangeRed}{HTML}{FF4500}
\definecolor{ImpDarkGreen}{HTML}{006400}
\definecolor{ImpSlateGray}{HTML}{708090}
\definecolor{ImpYellow}{HTML}{FFFF00}
\definecolor{ImpTurquoise}{HTML}{40E0D0}
\definecolor{ImpViolet}{HTML}{EE82EE}
\definecolor{ImpMediumSlateBlue}{HTML}{7B68EE}
\definecolor{ImpRed}{HTML}{FF0000}
\definecolor{ImpDarkOrange}{HTML}{FF8C00}
\definecolor{ImpSpringGreen}{HTML}{00FF7F}
\definecolor{ImpWhiteSmoke}{HTML}{F5F5F5}
\definecolor{ImpDeepSkyBlue}{HTML}{00BFFF}
\definecolor{ImpKhaki}{HTML}{F0E68C}
\definecolor{ImpPaleTurquoise}{HTML}{AFEEEE}
\definecolor{ImpLightPink}{HTML}{FFB6C1}
\definecolor{ImpLavender}{HTML}{E6E6FA}
\definecolor{ImpSalmon}{HTML}{FA8072}
\definecolor{ImpOrange}{HTML}{FFA500}
\definecolor{ImpPaleGreen}{HTML}{98FB98}

%----------------------------------------------------------------------------------------
%	HEADINGS IN IMPERIAL BLUE
%----------------------------------------------------------------------------------------

% Chapter titles use the class's own hooks (see MastersDoctoralThesis.cls).
\RenewDocumentCommand{\chapterfont}{}{\Huge\bfseries\color{ImperialBlue}}
\RenewDocumentCommand{\chapterprefixfont}{}{\LARGE\bfseries\color{ImperialBlue}}

% Section levels: replicate the book-class defaults and add \color, so spacing is
% untouched and no sectioning package (titlesec/sectsty) is needed.
\makeatletter
\renewcommand\section{\@startsection{section}{1}{\z@}%
  {-3.5ex \@plus -1ex \@minus -.2ex}{2.3ex \@plus.2ex}%
  {\normalfont\Large\bfseries\color{ImperialBlue}}}
\renewcommand\subsection{\@startsection{subsection}{2}{\z@}%
  {-3.25ex\@plus -1ex \@minus -.2ex}{1.5ex \@plus .2ex}%
  {\normalfont\large\bfseries\color{ImperialBlue}}}
\renewcommand\subsubsection{\@startsection{subsubsection}{3}{\z@}%
  {-3.25ex\@plus -1ex \@minus -.2ex}{1.5ex \@plus .2ex}%
  {\normalfont\normalsize\bfseries\color{ImperialBlue}}}
\makeatother

% Title-page rules in Imperial Blue.
\renewcommand{\HRule}{\textcolor{ImperialBlue}{\rule{\linewidth}{0.6pt}}}
\renewcommand{\decoRule}{\textcolor{ImperialBlue}{\rule{.8\textwidth}{.4pt}}}

%----------------------------------------------------------------------------------------
%	HYPERLINK COLOURS (override the class defaults at end of preamble)
%----------------------------------------------------------------------------------------
% hyperref has a single knob for internal cross-references: linkcolor covers figure, table,
% section, equation and table-of-contents links alike, so making figure/table numbers black
% necessarily makes every internal reference black. Citations and URLs keep their colour.
\AtBeginDocument{\hypersetup{
  linkcolor=black,          % internal refs (figures, tables, sections, equations, ToC)
  citecolor=ImperialNavy,   % citations
  urlcolor=ImperialBlue,    % URLs
}}
